% !TeX root = part7.tex
\documentclass[../final.tex]{subfiles}
\begin{document}
\enlargethispage{-2\baselineskip}
% Условия: Семинар_7.pdf, 24.10.26.
% Свой конспект: Desktop/4.jpg--6.jpg, верх 7.jpg.
% Второй конспект: Desktop/11.jpg,22.jpg,33.jpg,44.jpg,55.jpg.
% Нормировка плотности состояний согласована с двумя ветвями.
% В задаче 2 восстановлено инфракрасное ограничение конечного образца.
\practice{7}{24.10.2026}{Фононы двумерной решётки}

\rtheory{От колебаний атомов к фононам}
Полагаем $k_B=1$. Пусть $N$ одинаковых атомов массы $m$ образуют
решётку площади $S$, $n_2=N/S$. Рассматриваем две компоненты
смещения \emph{в плоскости} поверхности. Это две акустические ветви:
продольная $L$ и поперечная $T$. Колебания по нормали предполагаем
подавленными связью с подложкой. Такая модель соответствует
использованному в конспектах счёту $2N$ мод.

\subsection*{Нормальные координаты}
В гармоническом приближении потенциал раскладывается около
равновесия до второго порядка по смещениям. После диагонализации
квадратичной формы гамильтониан становится суммой осцилляторов:
\[
 H=\sum_\nu\left(\frac{P_\nu^2}{2m}
 +\frac{m\omega_\nu^2Q_\nu^2}{2}\right),\qquad
 E_{\nu,n}=\hbar\omega_\nu\left(n+\frac12\right).
\]
Фонон --- квант нормального колебания. Число фононов не сохраняется,
поэтому $\mu_{\mathrm{ph}}=0$ и
\[
 \overline n_\nu=\frac1{e^{\hbar\omega_\nu/T}-1},\qquad
 \overline E_\nu=\hbar\omega_\nu\left(\frac12+\overline n_\nu\right).
\]

\task[1]{Теплоёмкость двумерной решётки}
\condition
Найти вклад колебаний атомов двумерной решётки в теплоёмкость;
атомы адсорбированы на поверхности.

\subsection*{Двумерная модель Дебая}
Для ветви $s=L,T$ принимаем $\omega_s=c_sk$. Число состояний
в кольце радиуса $k$ и толщины $dk$ равно
\[
 d\mathcal N_s=\frac{S}{(2\pi)^2}\,2\pi k\,dk,
 \qquad D_s(\omega)=\frac{S\omega}{2\pi c_s^2}.
\]
На каждую ветвь приходится $N$ мод. Дебаевский радиус и частоты
отсечения определяются нормировкой
\[
 \frac{S k_D^2}{4\pi}=N,\qquad
 k_D=\sqrt{4\pi n_2},\qquad
 \omega_{Ds}=c_s k_D,\qquad\theta_s=\hbar\omega_{Ds}.
\]
Таким образом,
\[
 D_s(\omega)=\frac{2N\omega}{\omega_{Ds}^2},\qquad
 \int_0^{\omega_{Ds}}D_s(\omega)\,d\omega=N.
\]
Нулевая энергия $E_0=\sum_sN\theta_s/3$ от температуры не зависит
и не вносит вклада в теплоёмкость. Тепловая часть энергии:
\[
 E_{\mathrm{th}}=
 \sum_{s=L,T}\frac{2NT^3}{\theta_s^2}
 \int_0^{\theta_s/T}\frac{x^2\,dx}{e^x-1}.
\]
Дифференцирование среднего числа заполнения даёт удобную формулу
для всех температур в модели Дебая:
\[
 C_V=\sum_{s=L,T}2N\left(\frac{T}{\theta_s}\right)^2
 \int_0^{\theta_s/T}\frac{x^3e^x}{(e^x-1)^2}\,dx.
\]

\subsection*{Низкая и высокая температуры}
При $T\ll\theta_L,\theta_T$ заменяем верхние пределы бесконечностью.
Используя $\int_0^\infty x^2/(e^x-1)\,dx=2\zeta(3)$, находим
\[
 E_{\mathrm{th}}=4N\zeta(3)T^3\sum_{s=L,T}\theta_s^{-2},
 \qquad
 \boxed{C_V=12N\zeta(3)T^2\left(\theta_L^{-2}+\theta_T^{-2}\right)}.
\]
В отличие от трёхмерного закона Дебая $T^3$, здесь теплоёмкость
пропорциональна $T^2$, поскольку $D(\omega)\propto\omega$.
При $T\gg\theta_L,\theta_T$ каждая из $2N$ мод даёт вклад единица:
\[
 \boxed{C_V\longrightarrow2N}.
\]
Если, как в конспекте, положить $c_L=c_T=c_s$,
то $\theta_L=\theta_T=\theta_D$, и
\[
 D(\omega)=\frac{4N\omega}{\omega_D^2},\qquad
 \boxed{C_V=24N\zeta(3)\left(\frac{T}{\theta_D}\right)^2
 \quad(T\ll\theta_D)}.
\]
В высокотемпературном пределе более точно
\[
 C_V=2N\left[1-\frac1{24}\left(\frac{\theta_D}{T}\right)^2
 +O\left(\frac{\theta_D^4}{T^4}\right)\right].
\]
Переходя к кельвинам, следует положить
$\Theta_s=\hbar\omega_{Ds}/k_B$ и умножить формулы теплоёмкости
на $k_B$ при сохранении отношения $T/\Theta_s$.

\clearpage
\task[2]{Среднеквадратичное смещение атома}
\condition
Найти среднеквадратичное смещение атома из положения равновесия
в той же модели двумерной решётки.

\subsection*{Вклад одной моды и суммирование}
Для квантового осциллятора средние кинетическая и потенциальная
энергии равны. Поэтому
\[
 \frac{m\omega_\nu^2}{2}\langle Q_\nu^2\rangle
 =\frac12\overline E_\nu,\qquad
 \langle Q_\nu^2\rangle
 =\frac{\hbar}{2m\omega_\nu}
 \coth\left(\frac{\hbar\omega_\nu}{2T}\right).
\]
Нормируем собственные векторы условием
$\sum_i|\boldsymbol e_{i\nu}|^2=1$ и пишем
$\boldsymbol u_i=\sum_\nu\boldsymbol e_{i\nu}Q_\nu$.
Некоррелированность нормальных мод и эквивалентность атомов дают
\[
 \langle u^2\rangle\equiv\langle|\boldsymbol u_i|^2\rangle
 =\frac1N\sum_\nu\frac{\hbar}{2m\omega_\nu}
 \coth\left(\frac{\hbar\omega_\nu}{2T}\right).
\]
Это сумма двух компонент смещения, а не одна компонента.
Для изотропной решётки $\langle u_x^2\rangle=\langle u^2\rangle/2$.

\subsection*{Конечный образец}
Общее перемещение решётки исключаем фиксацией её центра масс.
В образце линейного размера $L$ минимальное ненулевое волновое
число имеет порядок $k_{\min}\sim2\pi/L$ (для периодических границ).
Соответственно $\omega_{\min,s}=c_sk_{\min}$.
Заменяя сумму интегралом, получаем
\[
 \langle u^2\rangle=
 \sum_{s=L,T}\frac{\hbar}{m\omega_{Ds}^2}
 \int_{\omega_{\min,s}}^{\omega_{Ds}}
 \coth\left(\frac{\hbar\omega}{2T}\right)\,d\omega.
\]
Интеграл можно взять явно. При $T>0$
\[
 \boxed{\langle u^2\rangle=
 \sum_{s=L,T}\frac{2T}{m\omega_{Ds}^2}
 \ln\frac{\sinh(\theta_s/(2T))}
 {\sinh(\hbar\omega_{\min,s}/(2T))}}.
\]
Формула учитывает и нулевые, и тепловые колебания.

\subsection*{Нулевые колебания}
В пределе $T\to0$ для конечного образца
\[
 \langle u^2\rangle_0=
 \sum_{s=L,T}\frac{\hbar}{m\omega_{Ds}^2}
 (\omega_{Ds}-\omega_{\min,s})
 \simeq\sum_{s=L,T}\frac{\hbar}{m\omega_{Ds}}.
\]
Для одинаковых скоростей $\langle u^2\rangle_0\simeq2\hbar/(m\omega_D)$.
При нулевой температуре длинноволновая часть интеграла конечна.

\subsection*{Тепловые колебания и размер образца}
В классическом пределе $T\gg\theta_s$ имеем
$\coth(\hbar\omega/(2T))\simeq2T/(\hbar\omega)$. Отсюда
\[
 \boxed{\langle u^2\rangle\simeq
 \sum_{s=L,T}\frac{2T}{m\omega_{Ds}^2}
 \ln\frac{k_D}{k_{\min}}
 =\frac{T}{2\pi m n_2}
 \left(\frac1{c_L^2}+\frac1{c_T^2}\right)
 \ln\frac{k_D}{k_{\min}}}.
\]
Если скорости одинаковы, это
\[
 \langle u^2\rangle\simeq\frac{4T}{m\omega_D^2}
 \ln\frac{k_DL}{2\pi}.
\]
При любой фиксированной $T>0$ низкочастотная часть остаётся
классической: интеграл содержит $d\omega/\omega$.
Поэтому в идеальной двумерной акустической модели
$\langle u^2\rangle$ растёт как $\ln L$ и при $L\to\infty$
расходится. Подстановка нуля в нижний предел при $T>0$
не даёт конечного ответа.

\remark[Границы модели.]
Если подложка закрепляет атомы также в плоскости, появляется
частотная щель, которая обрезает эту расходимость.
Если разрешена дополнительная ветвь колебаний по нормали,
её вклад нужно считать с её собственным законом дисперсии.
Начальная запись $3N$ координат в конспекте относится к общей
трёхмерной модели; для выбранной здесь модели движения в плоскости
число мод равно $2N$.
\end{document}
